% Generated draft: requires mathematical and source review.
\documentclass[11pt]{article}
\usepackage[margin=25mm]{geometry}
\usepackage{fontspec}
\setmainfont{Latin Modern Roman}
\usepackage{amsmath,amssymb,mathtools,unicode-math}
\setmathfont{Latin Modern Math}
\usepackage{enumitem,needspace,xurl,hyperref}
\setlength{\parindent}{0pt}
\setlength{\parskip}{5pt}
\setlength{\emergencystretch}{3em}
\newcommand{\N}{\mathbb{N}}
\newcommand{\Z}{\mathbb{Z}}
\newcommand{\Q}{\mathbb{Q}}
\newcommand{\R}{\mathbb{R}}
\newcommand{\C}{\mathbb{C}}
\newcommand{\F}{\mathbb{F}}
\newcommand{\A}{\mathbb{A}}
\newcommand{\PP}{\mathbb{P}}
\newcommand{\E}{\mathbb{E}}
\newcommand{\eps}{\varepsilon}
\newcommand{\dd}{\,\mathrm{d}}
\begin{document}
\section*{UnsolvedMath ID: 800192}
\textbf{Source problem code:} catalogue entry
\subsection{Definitions and mathematical statement}
For a precise working formulation, use the conventional irreducible holomorphic symplectic meaning of a hyperk\"ahler manifold: a compact, connected, simply connected K\"ahler complex manifold \(X\) such that
\[
H^0(X,\Omega_X^2)=\C\sigma,
\]
where \(\sigma\) is a holomorphic two-form that is nondegenerate at every point. This convention is not explicitly specified in the supplied excerpt; in particular, no projectivity hypothesis is imposed here.

For a K3 surface \(S\), write \(S^{[r]}\) for the Hilbert scheme of zero-dimensional subschemes of length \(r\), where \(r\geq 1\). For a complex torus \(T\) of complex dimension two, use the convention
\[
K_r(T)=\Sigma^{-1}(0),\qquad
\Sigma:T^{[r+1]}\longrightarrow T,
\]
where \(\Sigma\) sends a subscheme to the sum, with multiplicities, of its associated zero-cycle. The excerpt does not specify its indexing convention for \(K_r(T)\). Two compact complex manifolds are deformation equivalent if they can be connected by a chain of smooth proper holomorphic families over connected complex bases.

\textbf{Problem (Beauville, 2010, Question 1).} Find new examples of hyperk\"ahler manifolds. In light of the source's discussion of deformation types, the intended working target is to construct such an \(X\) whose deformation-equivalence class is different from those represented by \(S^{[r]}\), by \(K_r(T)\), and by the two O'Grady constructions mentioned in the source. The interpretation of ``new'' as a new deformation type is inferred from that discussion, rather than explicitly defined in the quoted question. The excerpt does not describe the two O'Grady constructions, so their precise identification remains a source dependency.

\textbf{Source qualifications.} The preceding passage concerns moduli spaces \(\mathcal M\) and their natural compactifications \(\overline{\mathcal M}\) arising in ``the two last examples,'' which are absent from the supplied excerpt. Within that context, the source states that a smooth \(\overline{\mathcal M}\) is hyperk\"ahler, citing [M], and is a deformation of \(S^{[r]}\) or \(K_r(T)\), attributing this to G\"ottsche--Huybrechts, O'Grady, Yoshioka, and others. It also states that a singular \(\overline{\mathcal M}\) admits a hyperk\"ahler desingularization only in the two cases considered by O'Grady, citing [KLS]. These assertions retain the source's unstated contextual hypotheses; they are not formulated here as claims about arbitrary moduli spaces or compactifications.
\subsection{Short English statement}
Beauville asks for new hyperk\"ahler manifolds. The surrounding discussion suggests seeking a new deformation type beyond the K3 Hilbert schemes, generalized Kummer examples, and the two cited O'Grady constructions. The precise meaning of ``new,'' the convention for hyperk\"ahler, and the hypotheses behind the moduli-space discussion require review against the complete source.
\subsection{Sources}
\begin{itemize}
\item Supplied local UnsolvedMath registry record, permanent ID \(800192\), exact source problem code \texttt{catalogue entry}, title ``Holomorphic Symplectic Geometry --- Question 1.'' Collection: AMR Open Problem Lists. The record attributes the problem to Beauville, with proposed year 2010, and supplies the status label ``open''; this label is not an independently verified current mathematical assessment.
\item Beauville, \emph{Holomorphic Symplectic Geometry: A Problem List} (2010), Question 1. Supplied source URL: \url{https://arxiv.org/abs/1002.4321}. The source record reports access on 2026-10-16 and extraction by \texttt{arxiv\_latex}; these are supplied provenance details, not a claim of independent source access.
\item The supplied passage cites [M] for the smooth compactification assertion and [KLS] for the restriction on hyperk\"ahler desingularizations. It names G\"ottsche--Huybrechts, O'Grady, Yoshioka, and others for the deformation assertion. Full bibliographic entries and the surrounding hypotheses are not supplied.
\end{itemize}
\end{document}
